Un amortiguador viscoso ejerce la fuerza \(F_d=-c x'\), con \(c>0\). Mediante excitaciones externas, que pueden ser distintas, se mantienen estas dos respuestas periódicas de un mismo sistema masa--resorte:
\[
x_1(t)=A\cos(\omega t),\qquad x_2(t)=A\cos(2\omega t),
\qquad A,\omega>0.
\]
No hay transitorio. La energía mecánica se repite al completar cada periodo y el trabajo de una fuerza es \(\int F x'\,dt\).

Compare la energía disipada por el amortiguador en los dos ensayos, primero durante el mismo intervalo \([0,2\pi/\omega]\) y después durante un periodo propio de cada respuesta. Explique por qué los dos cocientes no coinciden e indique el signo del trabajo del amortiguador y el trabajo neto que debe aportar la excitación externa en cada intervalo considerado.
